
\chapter*{Cómo usar esta sesión integrada}
\phantomsection\addcontentsline{toc}{chapter}{Cómo usar esta sesión integrada}\label{mod:ruta}
\begin{idea}[Pregunta práctica que justifica toda la sesión]
Un equipo técnico redacta bitácoras en español. Queremos saber si un modelo entrenado
sobre lenguaje general representa y predice esos textos con la misma eficacia que uno
adaptado al vocabulario del dominio. Antes de decir que un modelo es ``mejor'', debemos
\emph{definir unidades, aprender probabilidades, evaluar documentos no vistos y medir incertidumbre}.
La pregunta nació de un problema observable, no de una ecuación presentada sin contexto.
\end{idea}
\begin{center}
\begin{tikzpicture}[>=Stealth, node distance=6mm and 8mm,
  paso/.style={draw=azul, rounded corners,fill=azulclaro,minimum height=12mm,
  text width=31mm,align=center,font=\small}]
\node[paso](uno){Informes en español};
\node[paso,right=of uno](dos){Caracteres, tokens y pares};
\node[paso,right=of dos](tres){Conteos y probabilidades};
\node[paso,below=of tres](cuatro){BPC y diferencias por documento};
\node[paso,left=of cuatro](cinco){Informe con incertidumbre};
\node[paso,left=of cinco](seis){Pregunta de investigación};
\draw[->,thick](uno)--(dos);\draw[->,thick](dos)--(tres);
\draw[->,thick](tres)--(cuatro);\draw[->,thick](cuatro)--(cinco);
\draw[->,thick](cinco)--(seis);
\end{tikzpicture}
\end{center}
\noindent\textbf{Dos preguntas, dos mediciones:} los bigramas de caracteres
cuantifican una perdida predictiva (BPC), mientras que el
\hyperref[lab:ic10]{microexperimento BPE} cuantifica tokens por palabra
al cambiar el tokenizador. La segunda magnitud no mide por si sola calidad
probabilistica; separar ambas evita una conclusion indebida sobre IC10.

\noindent\textbf{Qué contiene este libro.} Este archivo \texttt{s01\_integrada.tex} es
\emph{la puerta de entrada} y el índice común de materiales que también compilan
por separado. Se puede leer únicamente la guía inicial, entrar después al
\hyperref[lab:mat]{laboratorio matemático}, ejecutar el
\hyperref[lab:comp]{laboratorio computacional} y su
\hyperref[lab:ic10]{microexperimento de tokenizacion}, consultar el
\hyperref[mod:sup]{suplemento} o resolver el
\hyperref[mod:taller]{taller}. La \hyperref[mod:invest]{ruta de investigación}
relaciona esos resultados con el documento de investigación y el objetivo de una
futura contribución científica en una revista de primer cuartil (Q1).

\begin{center}\small
\begin{tabular}{p{.23\textwidth}p{.65\textwidth}}
\toprule
\textbf{Componente} & \textbf{Función y evidencia que produce}\\\midrule
Guía y ejemplos & Construir el lenguaje matemático y comprender su utilidad.\\
Laboratorio matemático & Repetir cálculos a mano y demostrar identidades.\\
Laboratorio computacional & Generar resultados de ejemplo y archivos de auditoría en R y Python.\\
Taller + reporte & Interpretar cifras y producir un informe breve reproducible.\\
Suplemento & Justificar hipótesis, propiedades y límites de validez.\\
Investigación & Elegir dos áreas; identificar el camino a una pregunta novedosa y evaluable.\\\bottomrule
\end{tabular}
\end{center}
\begin{contraste}[Cuatro alcances que nunca se deben confundir]
\textbf{Cálculo demostrado:} se deduce de hipótesis declaradas.
\textbf{ILUSTRATIVO:} corpus o resultado fabricado para aprender.
\textbf{REAL REPRODUCIDO:} script ejecutado sobre un conjunto identificado
(que podría seguir siendo sintético).
\textbf{EVIDENCIA DE INVESTIGACIÓN:} requiere datos de dominio,
baselines fuertes, protocolo y validación externa. Los experimentos de esta
sesión pertenecen al segundo y tercer tipo; \textbf{no prueban la hipótesis de IC10}.
\end{contraste}
